\section{Auxiliary polynomials and local admissibility}\label{sec:local}

Throughout this section, $h\in\Z[x]$ is prime-intersective, of
degree $k\ge2$ and leading coefficient $a>0$.
Restricting a set to an arithmetic progression changes the polynomial
describing its forbidden differences.  We first choose progression steps
for which the transformed polynomials belong to a common auxiliary family.
We then select local polynomial values whose
arguments satisfy the unit condition required of primes.  These values will
supply the increments of the residue-label cycles used later.

\subsection{A family compatible with progression restriction}

We use the auxiliary-polynomial construction of Lucier~\cite{Lucier2006}
in its prime-intersective form in Rice~\cite[Section~2.1]{Rice2013}.
For every prime $p$, prime-intersectivity gives a root of $h$ in
$\Z_p^\times$, the units of the $p$-adic integers.  Indeed, the unit
roots modulo $p^j$ form a finitely branching tree under reduction, with
vertices at every depth.  Successively choosing a vertex that has descendants
at arbitrarily large depths gives a compatible sequence of roots.  Fix its
limit $\mathfrak z_p\in\Z_p^\times$, and write
\begin{equation}\label{eq:root-expansion}
 h(\mathfrak z_p+X)=\sum_{j=m_p}^k b_{j,p}X^j,
 \qquad b_{m_p,p}\ne0,\qquad 1\le m_p\le k.
\end{equation}
The coefficients lie in $\Z_p$.  Define a completely multiplicative
function $\lambda:\N\to\N$ by $\lambda(p)=p^{m_p}$.
For $l\ge1$, the Chinese remainder theorem specifies a unique integer
$r_l\in(-l,0]$ satisfying
\[
 r_l\equiv\mathfrak z_p\pmod{p^{v_p(l)}}\qquad(p\mid l).
\]
Here $v_p$ denotes the $p$-adic valuation.  In particular $(r_l,l)=1$;
for $l=1$ the convention is $r_1=0$.  Put
\begin{equation}\label{eq:auxiliary-definition}
 T_l(x)=r_l+lx,\qquad
 h_l(x)=\frac{h(T_l(x))}{\lambda(l)},\qquad
 a_l=\frac{a l^k}{\lambda(l)}.
\end{equation}

\begin{lemma}[Auxiliary family]\label{lem:auxiliary}
For every $l\ge1$, the polynomial $h_l$ belongs to $\Z[x]$, has
degree $k$ and positive leading coefficient $a_l$, and all its coefficients
are $O_h(a_l)$.  We have $h_1=h$ and $T_1(x)=x$.
For every $D\ge1$, there is an integer $i\in\{0,\ldots,D-1\}$ such that
\begin{equation}\label{eq:auxiliary-change}
 \begin{split}
 T_l(Dx-i)&=T_{lD}(x),\\
 h_l(Dx-i)&=\lambda(D)h_{lD}(x),\qquad
 D\mid\lambda(D),\quad \lambda(D)\le D^k.
 \end{split}
\end{equation}
Moreover, $b_{m_p,p}$ is a $p$-adic unit for all sufficiently large $p$,
with a threshold depending only on $h$.
\end{lemma}

\begin{proof}
To check integrality at $p\mid l$, write $u=v_p(l)$ and make the affine
change of variable
\[
 Y=\frac{r_l-\mathfrak z_p}{p^u}+\frac{l}{p^u}x.
\]
Its coefficients lie in $\Z_p$, and its slope is a unit.  Since
$v_p(\lambda(l))=um_p$, \eqref{eq:root-expansion} gives
\begin{equation}\label{eq:auxiliary-padic}
 h_l(x)=
 \left(\frac{\lambda(l)}{p^{um_p}}\right)^{-1}
 \sum_{j=m_p}^k b_{j,p}p^{u(j-m_p)}Y^j.
\end{equation}
Thus $h_l$ is integral at $p$.  If $p\nmid l$, its denominator
$\lambda(l)$ is a $p$-adic unit, so integrality is immediate.  This proves
$h_l\in\Z[x]$.  Its leading coefficient is $a_l$, and expanding
$h(r_l+lx)$ with $|r_l|<l$ bounds each coefficient by
$O_h(l^k/\lambda(l))=O_h(a_l)$.

The residues $r_l$ and $r_{lD}$ agree modulo $l$.  Their specified ranges
therefore give $r_{lD}=r_l-il$ for some $0\le i<D$.  Substitution and
complete multiplicativity prove \eqref{eq:auxiliary-change}; its last two
assertions follow from $1\le m_p\le k$.

For the final assertion, factor
$h=c\prod_i g_i^{d_i}$ over $\Z[x]$, where $c\ne0$ is an integer
and the $g_i$ are distinct primitive irreducible nonconstant polynomials.
Exclude primes dividing $c$ or the nonzero integer right-hand sides of
cleared B\'ezout identities for $g_i,g_j$ with $i\ne j$, and for
$g_i,g_i'$.  At a root $\mathfrak z_p$ for any remaining prime, exactly one
factor $g_i$ vanishes, its derivative is a unit, and all the other factors
have unit values.  Consequently $m_p=d_i$, and the coefficient of
$X^{m_p}$ in \eqref{eq:root-expansion} is a unit.
\end{proof}

Define the forbidden values at parameter $l$ by
\begin{equation}\label{eq:forbidden-family}
 \mathcal D_l=
 \{h_l(m):m\in\Z,\ m\ge1,\ T_l(m)\text{ is prime}\}.
\end{equation}
A set $A\subseteq[N]$ is \emph{avoiding at $l$} if
$(A-A)\cap\mathcal D_l\subseteq\{0\}$.  This condition passes to
progression coordinates: for every $D,n\ge1$ and every integer $c$, the set
\[
 A'=\{x\in[n]:\lambda(D)x+c\in A\}
\]
is avoiding at $lD$ whenever $A$ is avoiding at $l$.  In fact, a nonzero
difference $h_{lD}(m)$ in $A'$ would give the nonzero difference
$\lambda(D)h_{lD}(m)=h_l(Dm-i)$ in $A$.  Formula
\eqref{eq:auxiliary-change} preserves its prime argument, and
$Dm-i\ge1$ for $m\ge1$.  This also explains why no restriction at a zero
polynomial value is needed.

\subsection{Uniform cancellation at large primes}

We next estimate exponential sums over arguments $x$ satisfying
$p\nmid T_l(x)$, and over the subset on which $h_l'(x)$ is nonzero modulo
$p$.  Besides controlling local
Fourier coefficients, the latter estimates will produce zero-sum lists of
polynomial values: we seek one length $L$, independent of $l$ and $p$,
for which $L$ selected values sum to zero in each local residue ring.
At large primes cancellation supplies these lists; at the finitely many
smaller primes we will repeat a single value.  The selected values must
also come from argument classes on which values lift uniformly to higher
prime powers.

Fix
\begin{equation}\label{eq:local-exponents}
 \delta=\frac{1}{100k2^k},\qquad \gamma=\frac{\delta}{10}.
\end{equation}
For the remainder of this section, $l\ge1$ is arbitrary and $F=h_l$.
At all sufficiently large primes, $F\bmod p$ is nonconstant, uniformly
in $l$.  For $p\nmid l$, this follows from $p\nmid a$ and
\eqref{eq:auxiliary-definition}; for $p\mid l$, formula
\eqref{eq:auxiliary-padic} reduces to a unit multiple of $Y^{m_p}$.
When $p>k$, it follows that $F'\bmod p$ is not identically zero.
Call an argument $x\in\F_p$ \emph{allowed regular} if
\[
 p\nmid T_l(x)\qquad\text{and}\qquad F'(x)\ne0\pmod p.
\]
The unit condition excludes at most one residue: if $p\mid l$, it
excludes none because $(r_l,l)=1$.

\begin{lemma}[Local exponential sums]\label{lem:local-sums}
There is an integer $P_1>\max(k,2)$, depending only on $h$, such that for
every $l\ge1$ and every prime $p\ge P_1$, at least $p/2$ arguments in
$\F_p$ are allowed regular.  For every integer $v$ with $p\nmid v$,
\begin{equation}\label{eq:local-sums}
 \begin{split}
 \left|\mathbb E_{\substack{x\bmod p^j\\p\nmid T_l(x)}}
 e\left(\frac{vF(x)}{p^j}\right)\right|
 &\le p^{-8\delta j} \qquad(j\ge1),\\
 \left|\mathbb E_{\substack{x\bmod p\\x\text{ allowed regular}}}
 e\left(\frac{vF(x)}p\right)\right|
 &\le p^{-4\delta}.
 \end{split}
\end{equation}
Each expectation is normalized by the number of arguments in its stated
averaging set.
\end{lemma}

\begin{proof}
Take the threshold large enough for the preceding nonconstancy statements.
The derivative excludes at most $k-1$ residues, so at most $k$ arguments
fail to be allowed regular.  This gives the asserted cardinality after
increasing the threshold.

We first work modulo $p$.  If $k'$ is the degree of $F\bmod p$, then
the complete mean $S=\mathbb E_{x\bmod p}e(vF(x)/p)$ vanishes when
$k'=1$.  For $k'\ge2$, repeated differencing gives
\[
 |S|^{2^{k'-1}}
 \le\mathbb E_{u_1,\ldots,u_{k'-1}\bmod p}
 \left|\mathbb E_{x\bmod p}
 e\left(\frac{v\Delta_{u_1}\cdots\Delta_{u_{k'-1}}F(x)}p\right)\right|
 \le\frac{k'-1}{p},
\]
where $\Delta_u F(x)=F(x+u)-F(x)$.  The first inequality follows by
squaring, grouping the two arguments by their difference, and iterating
the triangle and Cauchy--Schwarz inequalities.  For the second, the final
phase is linear, with nonzero coefficient unless one of the $u_i$ is
zero: its coefficient is $v k'!$ times the leading coefficient of
$F\bmod p$ times $u_1\cdots u_{k'-1}$.  A union bound gives the displayed
probability.  Removing at most $k$ arguments changes a normalized mean
by $O_k(p^{-1})$.  Since $8\delta<2^{-(k-1)}$, this proves both assertions
of \eqref{eq:local-sums} at conductor $p$, for sufficiently large $p$.

For $j\ge2$, set $n=\lfloor j/2\rfloor$ and write
$x=x_0+p^{j-n}y$, with $x_0$ modulo $p^{j-n}$ and $y$ modulo $p^n$.
Taylor expansion over the integers gives
\[
 F(x_0+p^{j-n}y)\equiv F(x_0)+p^{j-n}yF'(x_0)\pmod{p^j}.
\]
The condition $p\nmid T_l(x)$ depends only on $x_0$.  Averaging over $y$
therefore annihilates every term for which $p^n\nmid F'(x_0)$.
We bound the proportion of the remaining $x_0$ by interpolation.

Put $b=\lceil n/(k-1)\rceil$.  The solutions of
$F'(x_0)\equiv0\pmod{p^n}$ occupy at most $k-1$ residue classes modulo
$p^b$.  Otherwise, choose $k$ integer solutions $x_1,\ldots,x_k$ in
distinct such classes.  In the Lagrange interpolation formula for the
degree-at-most-$(k-1)$ polynomial $F'$, every denominator
$\prod_{s\ne r}(x_r-x_s)$ has valuation at most
\[
 (k-1)(b-1)<n.
\]
Since every $F'(x_r)$ is divisible by $p^n$, each interpolating summand
has all coefficients divisible by $p$.  This contradicts
$F'\not\equiv0\pmod p$.

Since $b\le n\le j-n$, the class count applies to the full range of
$x_0$ modulo $p^{j-n}$.  The unit condition retains at least a proportion
$1/2$ of all arguments.
Consequently the conditional proportion that can survive $y$-averaging
is at most
\[
 2(k-1)p^{-b}
 \le 2(k-1)p^{-j/(3(k-1))},
\]
where $\lfloor j/2\rfloor\ge j/3$ for $j\ge2$.  Because
$8\delta<1/(3(k-1))$, one further increase of $P_1$ makes this at most
$p^{-8\delta j}$ for every $j\ge2$.  All thresholds were independent of
$l$.
\end{proof}

\subsection{Small primes and uniform lifting}

The preceding cancellation estimate will be used only for $p\ge P_1$.
For the finitely many smaller primes, we instead select a residue class
on which polynomial values lift uniformly.  The crucial point is that
its modulus can be fixed independently of $l$.

\begin{lemma}[Uniform lifting classes]\label{lem:local-lifting}
For every prime $p<P_1$, there is an integer $S_p\ge0$, depending only
on $h$ and $p$, such that, for every $l\ge1$, some integer $x$ satisfies
\begin{equation}\label{eq:small-prime-argument}
 p\nmid T_l(x),\qquad v_p(h_l'(x))\le S_p.
\end{equation}
Set
\begin{equation}\label{eq:local-modulus-exponents}
 E_p=2S_p+1\quad(p<P_1),\qquad E_p=1\quad(p\ge P_1).
\end{equation}
For any prime $p$, let $x$ be an integer satisfying
\eqref{eq:small-prime-argument} when $p<P_1$, or representing an allowed
regular argument when $p\ge P_1$.  Write $s=v_p(h_l'(x))$.
Every integer in the class $x\pmod{p^{E_p-s}}$ satisfies the unit
condition, and its $h_l$-value is congruent to $h_l(x)$ modulo $p^{E_p}$.
For each $j>E_p$, the uniform distribution on this argument class
modulo $p^j$ pushes forward under $h_l$ to the uniform distribution on
\[
 \{y\bmod p^j:y\equiv h_l(x)\pmod{p^{E_p}}\}.
\]
\end{lemma}

\begin{proof}
Fix $p<P_1$.  For a nonconstant polynomial over $\Z_p$, the
minimum valuation of its nonconstant coefficients is preserved by an
affine substitution of unit slope.  One inequality follows by expanding
the substitution, and the other by applying its inverse.  If $p\nmid l$,
this observation applied to \eqref{eq:auxiliary-definition} bounds that
minimum for $F=h_l$ by a constant depending only on $h,p$.  If $p\mid l$,
use \eqref{eq:auxiliary-padic}: before the unit affine substitution, its
degree-$m_p$ coefficient is a unit multiple of $b_{m_p,p}$, whose
valuation is fixed.  The same uniform bound therefore holds in this
case.  Differentiation increases the minimum valuation by at most
$\max_{1\le j\le k}v_p(j)$, so the minimum coefficient valuation of
$F'$ has a uniform upper bound $C_p\ge0$, which we take to be an integer.

There are at least $k$ distinct integers in $[1,kp]$ satisfying
$p\nmid T_l(x)$.  Interpolate $F'$, of degree at most $k-1$, at any
$k$ of them.  The interpolation denominators are products of differences
of integers in this fixed interval.  Let $B_p\ge0$ be a uniform integer
upper bound for their $p$-adic valuations.  If all $k$ values of $F'$ had
valuation greater than $C_p+B_p$, Lagrange interpolation would force
every coefficient of $F'$ to have valuation greater than $C_p$, a
contradiction.  Thus \eqref{eq:small-prime-argument} holds with
$S_p=C_p+B_p$.

Now let $p$ be arbitrary and $x,s$ be as in the statement.  In either
case $E_p\ge2s+1$.  Taylor expansion with integer coefficients gives
\begin{equation}\label{eq:local-permutation}
 F(x+p^{E_p-s}y)=F(x)+p^{E_p}\big(c_1y+pC(y)\big),
 \qquad p\nmid c_1,\quad C\in\Z[y].
\end{equation}
Indeed $c_1=p^{-s}F'(x)$, and each term of degree at least two contains,
after division by $p^{E_p}$, a factor $p^{E_p-2s}$.
The polynomial $G(y)=c_1y+pC(y)$ permutes the residues modulo every
power of $p$: for integers $y,z$, the quotient
$(G(y)-G(z))/(y-z)$ is a unit when $y\ne z$, so
\[
 v_p(G(y)-G(z))=v_p(y-z).
\]
It follows that $G$ is injective, hence bijective, on every finite
residue ring.  Formula \eqref{eq:local-permutation} proves constancy
modulo $p^{E_p}$.  The argument class preserves the unit condition
because $E_p-s\ge1$.

Finally, uniform arguments in this class modulo $p^j$ correspond to
uniform $y$ modulo $p^{j-E_p+s}$.  Each residue of $y$ modulo
$p^{j-E_p}$ occurs exactly $p^s$ times.  Applying the permutation $G$
modulo $p^{j-E_p}$ therefore gives precisely the asserted uniform
distribution of polynomial values.
\end{proof}

Call a residue modulo $p^{E_p}$ \emph{admissible at $l$} if it equals
$h_l(x)$ for an argument satisfying \eqref{eq:small-prime-argument}
when $p<P_1$, or for an allowed regular argument when $p\ge P_1$.
The equality is taken modulo $p^{E_p}$.  Thus every admissible residue
has an argument class satisfying the unit condition and the uniform lifting property of
Lemma~\ref{lem:local-lifting}.  For a modulus that is a product of
distinct factors $p^{E_p}$, admissibility means admissibility at every
factor.

\begin{lemma}[A fixed admissible cycle length]\label{lem:admissible-cycle}
Fix an integer $L>2$ such that
\begin{equation}\label{eq:cycle-length}
 4\delta L>2,\qquad p^{E_p}\mid L\quad(p<P_1).
\end{equation}
For every $l\ge1$ and every prime $p$, there are $L$ admissible residues
$y_0,\ldots,y_{L-1}$ modulo $p^{E_p}$, with repetitions allowed, such
that
\[
 y_0+\cdots+y_{L-1}=0\pmod{p^{E_p}}.
\]
\end{lemma}

\begin{proof}
For $p<P_1$, an admissible residue exists by
Lemma~\ref{lem:local-lifting}.  Repeat it $L$ times and use
$p^{E_p}\mid L$.  For $p\ge P_1$, take independent uniform allowed
regular arguments $x_0,\ldots,x_{L-1}$ modulo $p$.  Character
orthogonality and \eqref{eq:local-sums} give
\[
 \left|p\,\mathbb P\left(\sum_{j=0}^{L-1}F(x_j)=0\pmod p\right)-1\right|
 \le (p-1)p^{-4\delta L}<1.
\]
The probability is therefore positive, yielding the required sequence
of admissible values.
\end{proof}

The terminology refers to the residue labels
$z_0=0$ and $z_j=y_0+\cdots+y_{j-1}$ for $1\le j\le L$:
we have $z_L=z_0$, and the increment on the indexed edge $j\to j+1$
is the admissible value $y_j$.  Distinct indices need not carry distinct
residue labels.

We now keep $P_1$, the exponents $E_p$, and $L$ fixed.  For any later
threshold $P>P_1$, apply Lemma~\ref{lem:admissible-cycle} at each $p<P$
and combine the lists by the Chinese remainder theorem.  This gives $L$
admissible residues modulo $\prod_{p<P}p^{E_p}$ with sum zero.  Thus a
larger prime cutoff never requires a new cycle length.
